\section{Related Work \& Background}
\label{sec:related}

\paragraph{Machine Unlearning for LLMs.}
Machine unlearning was originally studied to remove specific training examples from
model memory for privacy reasons~\citep{cao2015towards}.
For LLMs, the goal has shifted toward \emph{concept-level forgetting}: suppressing
entire knowledge domains (biosecurity, copyrighted content) rather than individual
data points~\citep{liu2024rethinking}.
Methods differ in their optimization objective.
\begin{itemize}[leftmargin=1.5em, itemsep=2pt]
\item\emph{Gradient Difference} (GradDiff)~\citep{yao2023large} maximizes loss on forget
data while minimizing deviation on retain data.
\item\emph{Representation Mismatch Unlearning} (RMU)~\citep{li2024wmdp} disrupts the
internal representations of forget-domain inputs by steering them toward those of
random noise.
\item\emph{Representation Noising} (RepNoise)~\citep{rosati2024representation} injects
structured noise into the representations of harmful content during fine-tuning.
\item\emph{reinforcement-based forgetting} (RR)
\item\emph{task-aware regularization} (TAR)
\item\emph{latent-space adversarial perturbation} (RMU-LAT)
\item\emph{ensemble lifelong management} (ELM) 
\item\emph{preference-based forgetting} (PB\&J)
\end{itemize}
These eight methods are collected in the LLM-GAT evaluation suite~\citep{llmgat2024},
which provides standardized checkpoints fine-tuned from \textsc{Llama-3-8B-Instruct}.

\paragraph{The WMDP Benchmark.}
WMDP~\citep{li2024wmdp} is a multiple-choice benchmark covering biosecurity,
cybersecurity, and chemical weapon domains, designed to measure machine unlearning.
Successful unlearning is indicated by accuracy falling toward the random baseline
(25\% for 4-choice, 50\% for True/False).
Prior work using WMDP focuses exclusively on generation-level accuracy;
our contribution is to complement this with a hidden-state perspective.

\paragraph{Probing Classifiers.}
Probing classifiers~\citep{alain2017understanding} are shallow models---typically
logistic regression---trained on intermediate neural representations to test whether
a property is linearly decodable from activations.
Probing has revealed that syntactic structure is encoded in lower
layers~\citep{hewitt2019structural}, while factual associations are concentrated in
mid-to-late layers~\citep{meng2022locating}.
Most relevant to our work, \citet{levi2025jailbreak} showed that compliance-relevant
directions extracted from jailbreak initializations are linearly decodable, and
\citet{zou2023representation} demonstrated that safety-relevant concepts can be
represented as linear directions in activation space.
We apply probing to the question of whether unlearning erases or merely suppresses
these directions.

\paragraph{Representation Engineering and Activation Steering.}
\citet{zou2023representation} introduced \emph{representation engineering}, showing
that model behavior can be reliably modified by reading from or writing to specific
linear directions in activation space.
Related work on activation steering~\citep{turner2023activation} demonstrates that
hidden-state interventions can substantially change model outputs.
Taken together, this body of work implies that behavioral suppression via gradient
fine-tuning may leave underlying representational geometry intact---a hypothesis our
experiments directly test.

\paragraph{Knowledge Localization and Editing.}
A parallel line of work uses causal tracing~\citep{meng2022locating} and
gradient-based editing~\citep{meng2022mass} to identify and modify factual
associations within specific MLP layers.
These findings suggest that knowledge is encoded in localized, modifiable structures.
Unlearning methods that do not explicitly target these structures may therefore
produce only superficial behavioral changes, leaving the encoded knowledge intact
for a sufficiently capable readout.
